\documentclass[10pt,a4paper]{amsart}
\usepackage[margin=19mm]{geometry}
\usepackage{amsmath,amssymb,mathtools}
\usepackage[scr=rsfs,cal=euler]{mathalfa}
\usepackage{xfrac}
\usepackage[unicode,hidelinks,bookmarks=false]{hyperref}
\newcommand{\ie}{i.e.\ }
\newcommand{\cf}{cf.\ }
\newcommand{\resp}{respectively\ }
\newcommand{\Subsection}{\subsection}
\providecommand{\nobreakdash}{\nobreak\hbox{-}}
\setlength{\emergencystretch}{2em}
\multlinegap0pt
\usepackage{bm}
\usepackage{algorithm}\usepackage{algpseudocode}
\usepackage{amssymb}
\usepackage[normalem]{ulem}
\newtheorem{lemma}{Lemma}
\newtheorem{proposition}{Proposition}
\title{Heavy Tailed Homogeneous Structural Causal Models}
\author{Vishal Routh, Shuyang Bai}
\date{Source: arXiv:2604.04118v1 (CC BY 4.0)}
\begin{document}
\maketitle

\noindent\emph{Source and licence.} Heavy Tailed Homogeneous Structural Causal Models. Version arXiv:2604.04118v1 (CC BY 4.0), distributed by arXiv under the Creative Commons Attribution 4.0 International licence, http://creativecommons.org/licenses/by/4.0/, as recorded in the arXiv licence metadata for this exact version and shown on the abstract page https://arxiv.org/abs/2604.04118v1. Full licence notice: \texttt{licenses/arxiv-2604.04118-CC-BY-4.0.txt}.\par\smallskip

\noindent\textit{Editorial scope.} Counted unit: the article's proposition Pro:alg (source0.tex lines 473--475). The article's complete original proof of it is delivered with it (source0.tex lines 476--500, one \texttt{proof} environment of 25 lines). No mathematics is written, repaired or re-derived: the statement and the proof are the article's own lines, in the article's own notation, and no retained line is substituted or re-flowed.  Retained uncounted as the source-local context the proof needs: lines 432--447; lines 454--463.  Nothing was omitted from any retained span, so the package carries no omission record.  No retained line contains a citation, so the wrapper carries no bibliography and no bibliography entry.  No retained line contains a figure environment, an image-inclusion command or drawing code, so no image is delivered and the package declares no asset dependency.  Editorial formatting only, in addition: an AMS \texttt{amsart} preamble that restates the article's own declarations and macros used by the retained lines, namely the environments \texttt{lemma}, \texttt{proposition} on separate counters, together with the standard packages those lines need.  The retained environments print in this document as Lemma 1, Proposition 1 in order of appearance, whereas the article prints the same items with the numbering of its own full document.  The complete native source of the article as distributed in the arXiv e-print for this exact version is bundled unchanged as \texttt{sources/197860/source0.tex}.
 \begin{algorithm}[h]
 \caption{Find $\widetilde{F}$ from $\boldsymbol{\Gamma}^*$}
\label{alg:findF}
 \begin{algorithmic}[1]
 \For{$\nu = 0, \dots, d-1$}
     \For{each $j \in V$ with $|{an}(j)| = \nu$}
         \For{each $i \in V \setminus {De}(j)$}
             \State $\widetilde{F}_{ji} \gets 0$
         \EndFor
         \For{each $i \in {De}(j)$}
             \State $\widetilde{F}_{ji} \gets \Gamma^*_{ij} - \sum_{k \in {an}(j)} \widetilde{F}_{ki}$
         \EndFor
     \EndFor
 \EndFor
 \end{algorithmic}
 \end{algorithm}

\begin{lemma}\label{lem:ancestor_set_gamma}
For every $j \in V$,
\begin{equation}
an(j)= \{\, i \in V \setminus \{j\} : \Gamma^*_{ij}=1 \,\}.
\end{equation}
In particular,
\begin{equation}
|an(j)| = \sum_{i \in V \setminus \{j\}} \mathbf{1}\{\Gamma^*_{ij}=1\}.
\end{equation}
\end{lemma}

\begin{proposition}\label{Pro:alg}
Algorithm \ref{alg:findF} correctly recovers the standardized   AIR matrix $\widetilde{F}$    from the standardized CTC  matrix $\boldsymbol{\Gamma}^*$.
\end{proposition}

\begin{proof} 
    We proceed to prove the correctness of this algorithm by induction on the size of the ancestor set, $|{an}(j)|$ which is determined by  Lemma~\ref{lem:ancestor_set_gamma}
    
    $\bullet$  Base case. 
    
    Suppose $j \in V$ such that $|{an}(j)| = 0$ (i.e., $j$ is a root node). For $i \in V \setminus {De}(j)$, we trivially have $\widetilde{F}_{ji} = 0$. For $i \in {De}(j)$, since $j$ is an ancestor of $i$ and has no ancestors itself, the intersection simplifies to ${An}(i) \cap {An}(j) = {An}(j) = \{j\}$. Therefore, evaluating the causal tail coefficient yields:
    \[
        \Gamma^*_{ij} = \sum_{h \in {An}(i) \cap {An}(j)} \widetilde{F}_{hi} = \widetilde{F}_{ji}.
    \]
    Thus, $\widetilde{F}_{ji}$ is fully identified for all root nodes.
    
     $\bullet$   Inductive step.
     
     Assume that for all nodes $j \in V$ with $|{an}(j)| \le \nu$, the entries $\widetilde{F}_{ji}$ have been correctly obtained. Consider a node $j' \in V$ with $|{an}(j')| = \nu + 1$. 
    
    For $i \in V \setminus {De}(j')$, it trivially holds that $\widetilde{F}_{j'i} = 0$. For $i \in {De}(j')$, the intersection of ancestors is ${An}(i) \cap {An}(j') = {An}(j')$. Expanding the causal tail coefficient gives:
    \[
        \Gamma^*_{ij} = \sum_{h \in {An}(i) \cap {An}(j')} \widetilde{F}_{hi} = \sum_{h \in {An}(j')} \widetilde{F}_{hi} = \widetilde{F}_{j'i} + \sum_{h \in {an}(j')} \widetilde{F}_{hi}.
    \]
    Observe that if $h \in {an}(j')$, then $h$ is a strict ancestor of $j'$, which  implies $|{an}(h)| \le \nu$ due to the acyclic nature of the graph. By our inductive hypothesis, the values $\widetilde{F}_{hi}$ in the summation are already known. Therefore, $\widetilde{F}_{j'i}$ is uniquely and deterministically identified. This concludes the proof.



    
\end{proof}

\end{document}
